\chapter{立项依据}

\section{研究意义}

随着无人机感知、控制、计算和通信技术不断发展，无人机已逐渐由传统人工遥控飞行平台向具备自主感知、自主决策、自主规划和群体协同能力的智能无人系统发展。其应用范围也由早期航拍测绘、遥感监测等相对结构化任务，逐步扩展至低空巡检、应急救援、物流运输、城市治理、复杂区域搜索与探测等具有更强环境动态性和任务实时性的场景。在此类任务中，无人机通常难以获得完整且长期有效的环境先验信息，需要依赖机载传感器持续感知周围环境，并在有限机载计算资源条件下实时生成安全、平滑且满足动力学约束的飞行轨迹。因此，复杂环境下的自主运动规划已成为制约无人机智能化水平和实际应用能力的关键技术之一\cite{nie2025review,arshid2025review}。

无人机运动规划技术经历了由几何路径搜索向动力学轨迹规划，再向在线局部自主决策的发展过程。早期图搜索、人工势场和采样规划方法主要解决静态或具有较完整先验地图条件下的可行路径搜索问题；随着四旋翼高机动飞行需求提高，Minimum Snap、多项式轨迹优化等方法进一步考虑轨迹连续性和动力学可执行性。近年来，EGO-Planner、EVA-Planner、PE-Planner等实时局部规划方法显著增强了无人机在复杂环境中的在线规划能力，Primitive-Swarm进一步表明基于运动基元的轻量化方法能够兼顾规划实时性和大规模系统扩展能力\cite{egoPlanner,evaPlanner,pePlanner,primitiveSwarm}。

然而，实际无人机运行环境并非静态。人员、车辆、其他机器人以及其他无人机均可能构成动态障碍，其未来位置和运动方向随着时间持续变化。如果规划器仅依据当前障碍状态进行反应式避障，即使能够维持较高规划频率，也可能由于决策短视而发生提前量不足、往复规避甚至碰撞。因此，近年来动态环境规划逐渐由基于当前观测的避障向预测驱动规划发展。DPMPC-Planner利用预测时域处理动态障碍；SANDO面向动态未知环境研究安全自主轨迹规划；DYNUS进一步将不确定性感知引入动态未知环境；SPOT通过障碍威胁建模增强规划过程中的风险感知；HEPP则重点提升高速避障任务中的感知和规划效率\cite{dpmpc,sando,dynus,spot,hepp}。

动态障碍预测虽然能够提高规划前瞻性，却同时引入新的问题：预测本身并不可靠。传感器测量噪声、状态估计误差、有限历史观测、运动模型偏差以及障碍自主行为变化，均可能使预测轨迹与实际运动产生偏差。因此，“能够预测未来位置”并不意味着“能够安全利用预测结果”。规划系统还需要回答预测结果具有多大可信程度，以及当预测不确定性增加时是否应主动扩大安全裕度。Thomas等提出的时空占据栅格地图方法通过未来时空占据描述动态环境，为规划器提供区别于当前占据地图的未来状态表达；后续研究进一步将未来场景预测向时空风险表示拓展。Wu等则将时空占据栅格应用于多智能体动态环境轨迹规划，将动态障碍与其他智能体计划轨迹统一为未来时空占据信息\cite{thomas2022sogm,thomas2023future,wu2024sogm}。

另一方面，当无人机任务由单机扩展至无人机集群后，单机环境不确定性之外还会产生新的协同信息不确定性。多无人机分布式规划依赖邻机位置、速度、未来轨迹和任务状态等信息进行冲突判断及协同决策。EGO-Swarm和MADER等方法通过去中心化规划降低了对中央节点的依赖，但并没有消除通信依赖，而是将通信需求转化为局部信息共享需求\cite{egoSwarm,mader}。实际无线通信环境中，通信范围、遮挡、网络负载、链路质量和设备状态变化会造成时延、数据丢失及更新不及时。此时，无人机接收到的邻机状态虽然在数据层面仍“存在”，但可能已经无法准确代表邻机当前真实状态。因此，多机规划需要解决的问题不再只是“能否获得邻机信息”，而是“当前邻机信息是否仍具有规划价值”。

基于上述认识，本研究面向动态复杂环境中的两类信息可靠性问题展开：一方面研究动态障碍预测不确定性，使无人机能够根据障碍未来状态的可信程度调整局部避障决策；另一方面研究通信状态导致的协同信息有效性变化，使多无人机系统能够根据邻机共享信息的新鲜程度和可靠性动态调整协同行为。两个研究方向虽然分别面向环境侧和协同侧，但其共同本质均是由单纯利用状态信息，转向同时感知信息可靠程度，并根据信息可靠程度调整规划决策。

从理论角度看，本研究有助于推动无人机运动规划由确定性状态驱动向可靠性感知决策发展，并研究不确定性信息如何以较低计算代价进入实时局部规划。从工程应用角度看，研究成果有望提高无人机在动态障碍、感知偏差和非理想通信条件下的持续自主运行能力，为复杂低空环境中的巡检、搜索、救援和多无人机协同作业提供技术支撑。

\section{国内外研究现状分析及存在问题}

\subsection{无人机自主规划与动态未知环境避障研究现状}

\subsubsection{国外研究现状}

无人机自主规划最初主要建立在移动机器人路径规划理论基础之上。传统图搜索方法通过离散环境寻找从起点到终点的可行路径，人工势场方法依靠目标吸引和障碍排斥实现局部避障，RRT等采样方法则提升了高维连续状态空间中的规划能力。随着无人机飞行速度和机动性提高，研究进一步由几何路径转向满足动力学约束的连续轨迹生成。

进入实时自主规划阶段后，EGO-Planner通过避免显式构造ESDF降低局部优化开销，EVA-Planner进一步增强复杂环境适应能力，PE-Planner则在复杂动态环境中提升规划性能。近年Primitive-Swarm采用超轻量运动基元方法面向大规模空中集群实现快速规划，说明基于有限候选动作和局部评估的规划机制在兼顾实时性和扩展性方面具有明显优势\cite{egoPlanner,evaPlanner,pePlanner,primitiveSwarm}。

随着研究环境由静态向动态演进，障碍未来运动状态逐渐成为规划决策的重要信息。DPMPC-Planner采用模型预测思想处理动态障碍；Wakabayashi等进一步考虑障碍速度不确定性；SANDO、DYNUS、SPOT和HEPP等工作则从动态未知环境、不确定性感知、威胁建模和高速感知规划等方向进一步推进研究\cite{dpmpc,wakabayashi2023,sando,dynus,spot,hepp}。Thomas等提出的SOGM从环境表示角度将当前空间占据扩展为未来时空占据，Wu等进一步将其用于多智能体轨迹规划，形成“感知—预测—未来环境表示—轨迹规划”的技术链\cite{thomas2022sogm,wu2024sogm}。

\subsubsection{国内研究现状}

国内近年来也围绕三维路径规划、动态避障、集群避障和复杂环境规划进行了较多研究。聂虹宇等从环境感知和运动规划角度对多旋翼无人机研究进行了系统综述\cite{nie2025review}。殷泽阳等进一步面向行为未知动态障碍开展无人机编队路径规划研究，通过未来状态预测及可达包络描述障碍可能到达区域，将动态障碍运动状态不确定性显式纳入规划过程\cite{yin2025envelope}。

总体而言，国内外动态规划已经从静态路径搜索逐渐发展到在线局部规划和预测驱动规划。但是，未来障碍状态仍主要依赖有限观测和运动模型获得，预测误差将直接影响碰撞判断。因此，如何进一步表征预测可信程度，并将其有效用于实时规划，是动态未知环境研究需要继续解决的问题。

\subsection{不确定性感知与风险约束规划研究现状}

\subsubsection{国外研究现状}

不确定性规划的核心问题是：在系统无法准确获得未来环境状态时，如何保证轨迹仍具有足够安全裕度。机会约束方法为这一问题提供了重要理论基础。Gutow和Rogers将Koopman算子和机会约束运动基元规划结合，将数据驱动动态建模引入运动基元安全选择\cite{gutow2020koopman}。Nguyen等研究了基于深度碰撞预测的运动基元导航，使未来碰撞风险直接影响候选动作选择\cite{nguyen2022collision}；Liu等研究不确定环境下无人机运动规划中的紧致碰撞概率计算\cite{liu2023tight}；Wakabayashi等将动态障碍速度不确定性纳入机会约束\cite{wakabayashi2023}；Xia等针对动态不确定环境建立风险评估与MAV运动规划机制\cite{xia2024risk}；Multi-Uncertainty Aware Autonomous Cooperative Planning进一步考虑多来源不确定因素对自主协同规划的影响\cite{zhang2024multiuncertainty}。

从另一条技术路线看，未来环境预测正在与风险空间表示结合。SOGM利用时空占据表示未来动态环境状态，Thomas等后续研究进一步将未来场景预测转换为时空风险信息，使规划器能够根据不同空间位置和未来时间的风险程度进行路径选择\cite{thomas2022sogm,thomas2023future}。

\subsubsection{国内研究现状}

国内关于无人机预测不确定性的研究数量相对传统路径规划仍较少，但已经出现从未来单点预测向安全区域和不确定性边界建模发展的趋势。殷泽阳等通过预测动态障碍运动状态并传播其不确定性，计算预测时域内的动态可达包络，再将其作为编队路径规划中的安全约束\cite{yin2025envelope}。

现有不确定性感知和风险约束方法已经证明，将预测可信程度纳入规划能够提升复杂环境安全性，但部分方法依赖概率分布传播、机会约束及连续轨迹优化，在动态障碍数量增多时会增加实时计算负担；同时，不确定性建模和轨迹规划往往相对独立。为此，本研究拟以轻量运动基元局部规划为基础，将动态障碍预测不确定性直接转化为候选轨迹风险，同时利用TTC和最小安全距离构建硬安全层，在不采用高计算开销全局概率优化的条件下提高动态未知环境规划可靠性。

\subsection{多无人机协同规划研究现状}

\subsubsection{国外研究现状}

多无人机协同研究早期主要关注群体一致性、编队保持和群体运动控制。随着无人机自主决策能力提高，研究重点逐渐由“如何保持队形”扩展至“如何在复杂环境中独立规划同时保持整体安全”。EGO-Swarm在EGO-Planner实时局部规划基础上扩展至无人机集群，通过去中心化局部轨迹优化实现复杂障碍环境自主飞行；MADER则针对多智能体和动态环境建立去中心化轨迹规划框架，使不同机器人通过共享计划轨迹完成冲突避免\cite{egoSwarm,mader}。

进一步地，多智能体规划逐渐由“交换当前状态”向“预测未来交互”发展。Zhu等研究交互感知轨迹预测，使规划器能够提前估计其他智能体未来行为\cite{zhu2021interaction}；Wu等将其他机器人规划轨迹与动态障碍未来状态统一写入SOGM，并在统一时空占据表示下进行分布式规划\cite{wu2024sogm}。

\subsubsection{国内研究现状}

国内多无人机规划近年来重点关注集群避障、编队调整、协同航迹生成和复杂环境任务执行。穆凌霞等针对多无人机编队避障和编队重构问题开展研究，使无人机集群在障碍环境中能够根据环境约束调整编队结构\cite{mu2024formation}。徐蜜从分布式架构出发研究无人机集群协同搜索运动目标问题，并在局部通信条件下考虑信息共享和分布式协同决策\cite{xu2025search}。

总体来看，多无人机轨迹规划已经由集中式规划逐渐向去中心化、自主化方向发展，但分布式架构并没有消除无人机之间的信息依赖。当通信发生时延、丢包或短时中断时，不同无人机掌握的系统状态可能不再一致，导致规划器依据过期邻机轨迹进行冲突检测。因此，多无人机规划需要进一步发展为能够适应非理想信息交换条件的分布式规划。

\subsection{通信约束下多无人机协同规划研究现状}

\subsubsection{国外研究现状}

通信是多无人机分布式协同的重要基础。Le Ny等研究通信约束下无人系统自适应部署，说明通信网络结构能够直接影响群体协同行为\cite{leny2012communication}。Serra-Gomez等提出“与谁通信”的方法，通过选择具有更高任务价值的信息交互对象，提高有限通信资源条件下多机器人碰撞规避效果\cite{serragomez2020whom}。Mikkelsen等进一步将Fiedler值用于通信感知轨迹规划，使网络连通性直接进入轨迹优化约束\cite{mikkelsenCommunication}。

Guo等研究时变通信与能源约束下的多无人机协同路径规划\cite{guo2024communication}；Javaid等对协同无人机通信和控制的发展进行了系统总结\cite{javaid2023communication}；Mudassir Ejaz等进一步从跨层视角综述通信感知多无人机路径规划\cite{mudassir2026review}。另一方面，Robust MADER针对通信延迟条件下不同机器人获得的轨迹信息不同步问题，对MADER分布式规划机制进行增强，使规划系统能够在延迟条件下维持更高安全性\cite{robustMader}。CoCoPlan则进一步从协调和通信联合自适应角度研究动态未知环境中的多机器人系统\cite{cocoplan2026}。

\subsubsection{国内研究现状}

国内通信约束无人机研究已经开始从无线通信性能优化逐渐向通信与轨迹联合规划发展。刘炫麟围绕无人机辅助无线通信系统的轨迹规划开展研究，从无人机运动轨迹与无线通信性能耦合角度研究联合优化问题\cite{liu2025uavcomm}。胡昊南等在无线传感器网络多无人机轨迹优化问题中引入信息年龄（Age of Information，AoI）评价数据新鲜程度，并根据AoI开展多无人机轨迹优化\cite{hu2024aoi}。刘兴宇等提出基于通信功率自适应的无人机集群协同导航控制方法，通过通信能力调整影响无人机邻居关系和集群协同过程\cite{liu2024power}。徐蜜则针对分布式无人机集群协同搜索运动目标问题，在局部通信条件下研究信息共享、信息融合和协同决策\cite{xu2025search}。

总体来看，通信约束多无人机研究已经经历“通信连通保障—通信资源优化—通信感知规划—非理想通信鲁棒规划”的发展过程。然而，从运动规划角度看，通信性能与规划信息有效性之间并不能简单等价。通信链路保持连接，并不意味着无人机掌握的邻机状态一定足够新鲜；同样，瞬时链路质量下降也不意味着所有历史信息立即失效。因此，需要将通信状态进一步映射为规划层的信息有效性，并建立“通信状态—信息有效性—协同策略—轨迹规划”的完整决策链。

\subsection{国内外研究现状总结及项目切入点}

综合国内外研究现状，可以将动态复杂环境下无人机自主规划当前面临的问题归纳为两类。第一类是环境信息可靠性问题。动态规划已经由当前状态反应式避障发展至未来状态预测驱动规划，但动态障碍未来运动不可避免具有不确定性。现有机会约束、碰撞概率、风险评估、SOGM及时空风险表示等研究已经证明，利用预测不确定性能够提高规划安全性；然而，在资源受限无人机上，如何将不确定性以足够轻量的方式直接用于高频局部轨迹选择，仍需要进一步研究。

第二类是协同信息可靠性问题。多无人机已经由集中式控制逐步发展至分布式规划，但其协同决策仍依赖邻机状态及轨迹共享。通信时延、丢包和更新间隔会使信息逐渐偏离邻机真实状态，因此“拥有邻机信息”并不意味着“信息仍然可靠”。现有研究已分别从连通性、通信对象、AoI、通信功率和时延鲁棒等角度解决部分问题，但如何统一评价共享信息有效程度，并据此调整协同规划行为仍具有进一步研究空间。

因此，本项目从“环境预测可靠性”和“协同信息可靠性”两个层面切入：一是针对动态障碍建立短时运动预测与不确定性表征，将预测可信程度转化为候选运动基元的时空风险，并结合TTC和最小安全距离构建轻量化软硬融合安全规划机制；二是针对多无人机通信约束，建立通信状态与邻机信息有效性之间的映射关系，根据共享信息时效性和可靠程度动态调整协同强度、安全裕度以及邻机轨迹利用方式。

\section{应用方向及应用前景}

本项目面向复杂动态低空环境下无人机自主运行需求，研究成果可应用于城市低空巡检、应急救援、无人机集群搜索探测等场景。

在城市和园区低空巡检任务中，车辆、人员、其他无人机以及临时设施会形成持续变化的动态障碍。传统基于静态地图或单一确定性预测的规划方法难以长期适应此类环境。预测不确定性感知规划能够根据动态障碍未来行为的不确定程度主动调整避障裕度，从而提高复杂动态空间中的安全飞行能力。

在灾害救援和复杂区域搜索任务中，无人机通常无法获得完整先验地图，通信还可能受到建筑物、地形遮挡以及基础设施损坏影响。无人机既需要依赖机载感知完成局部自主避障，又需要通过有限通信共享任务和状态信息。因此，动态障碍不确定性规划和通信状态感知协同能够共同提高无人机集群在非理想环境中的持续任务能力。

在多无人机自主协同巡检、搜索、监测等任务中，随着无人机数量增多，固定通信和固定协同机制容易受到网络负载和通信范围限制。根据信息时效性和可靠程度动态调整协同方式，有望减小大规模集群对持续高质量通信的依赖，提高系统扩展性和容错能力。

本项目拟基于Gazebo、PX4和ROS2构建从环境感知、动态目标状态估计、预测、不确定性规划到多机通信协同的完整闭环验证体系。研究成果后续可进一步向真实无人机平台迁移，为动态复杂低空环境中的自主避障和多机协同提供算法及工程实现基础。
